\documentclass[11pt]{article}
\usepackage[margin=1in]{geometry}
\usepackage{xcolor}
\usepackage{tcolorbox}
\usepackage{hyperref}
\usepackage{enumitem}
\usepackage{booktabs}
\usepackage{array}
\usepackage{amsmath}

% Define OSU orange color
\definecolor{osaorange}{RGB}{220,115,38}

% Configure hyperref colors
\hypersetup{
    colorlinks=true,
    linkcolor=blue,
    urlcolor=blue,
    citecolor=blue
}

\title{}
\author{}
\date{}

\begin{document}

% Orange header box with black outline
\begin{tcolorbox}[colback=osaorange,colframe=black,boxrule=2pt,arc=0pt,left=6pt,right=6pt,top=10pt,bottom=10pt]
\centering
\textcolor{white}{\textbf{\Large Week 8 Assignment - ANSWER KEY}}\\\
\textcolor{white}{\textbf{\large Contingency Tables and Chi-Square Tests}}
\end{tcolorbox}

\vspace{8pt}

\noindent\textbf{Course:} H524 Introduction to Biostatistics\\
\textbf{Total Points:} 100 points

\vspace{8pt}

\section{Part A: Multiple Choice (70 points)}

\subsection*{Contingency Table Basics (10 points)}
A contingency table displays the relationship between:

\textbf{ANSWER: C - Two categorical variables}

\textit{Explanation:} Contingency tables are used to display the relationship between two categorical variables, showing the frequency distribution across categories.

\vspace{0.5cm}

\subsection*{Expected Frequencies (10 points)}
In a chi-square test, the expected frequency for a cell is calculated as:

\textbf{ANSWER: A - (Row total × Column total) / Grand total}

\textit{Explanation:} Expected frequency formula: $E_{ij} = \frac{(\text{Row } i \text{ total}) \times (\text{Column } j \text{ total})}{\text{Grand total}}$

\vspace{0.5cm}

\subsection*{Chi-Square Test Assumptions (10 points)}
Which of the following is a requirement for conducting a chi-square test?

\textbf{ANSWER: B - All expected cell frequencies should be at least 5}

\textit{Explanation:} The chi-square test assumes expected counts $\geq 5$ in all cells. If this is violated, use Fisher's exact test instead.

\vspace{0.5cm}

\subsection*{Degrees of Freedom (10 points)}
For a 3×4 contingency table (3 rows, 4 columns), what are the degrees of freedom for the chi-square test?

\textbf{ANSWER: B - 6}

\textit{Explanation:} $df = (r-1)(c-1) = (3-1)(4-1) = 2 \times 3 = 6$

\vspace{0.5cm}

\subsection*{Interpreting Chi-Square Results (10 points)}
A chi-square test yields $p = 0.032$. At $\alpha = 0.05$, what can we conclude?

\textbf{ANSWER: B - There is a significant association between the variables}

\textit{Explanation:} Since $p = 0.032 < 0.05$, we reject $H_0$ and conclude there is significant evidence of association.

\vspace{0.5cm}

\subsection*{2×2 Table Setup (10 points)}
In a case-control study examining smoking and lung cancer, which variable should form the rows of a 2×2 table?

\textbf{ANSWER: C - Either variable can form rows}

\textit{Explanation:} The orientation of a 2×2 table (which variable forms rows vs columns) does not affect the statistical test results.

\vspace{0.5cm}

\subsection*{Fisher's Exact Test (10 points)}
When should Fisher's exact test be used instead of the chi-square test?

\textbf{ANSWER: B - When expected cell frequencies are small (less than 5 in one or more cells)}

\textit{Explanation:} Fisher's exact test is appropriate when chi-square assumptions are violated (expected counts $< 5$).

\newpage

\section{Part B: R Computational Problems (30 points)}

\textbf{Note:} Some problems require students to use R to compute the answer.

\subsection*{Question 8: Fisher's Exact Test in R (15 points)}

A small clinical trial tests a new treatment for a rare disease with the following results:

\begin{center}
\begin{tabular}{l|cc}
& \textbf{Improved} & \textbf{Not Improved} \\
\hline
\textbf{Treatment} & 7 & 3 \\
\textbf{Control} & 2 & 8 \\
\end{tabular}
\end{center}

Use R to perform Fisher's exact test on this data. What is the two-sided p-value?

\begin{enumerate}[label=\Alph*.]
\item 0.0032
\item 0.0698 \textbf{(CORRECT)}
\item 0.1524
\item 0.5000
\end{enumerate}

\vspace{0.5cm}

\textbf{Expected R code solution:}
\begin{verbatim}
data <- matrix(c(7, 3, 2, 8), nrow=2, byrow=TRUE)
fisher.test(data)
\end{verbatim}

\textit{R output:}
\begin{verbatim}
	Fisher's Exact Test for Count Data

data:  data
p-value = 0.06978
alternative hypothesis: true odds ratio is not equal to 1
95 percent confidence interval:
   0.8821175 127.0558418
sample estimates:
odds ratio
  8.153063
\end{verbatim}

\textit{Explanation:} The two-sided p-value is 0.0698 (rounded to 0.0698). Since $p > 0.05$, we fail to reject $H_0$ at the 0.05 level. There is no statistically significant evidence of treatment effectiveness, though the p-value is close to 0.05 and the odds ratio suggests a strong effect. A larger study would be warranted.

\vspace{1cm}

\subsection*{Question 9: Chi-Square Test in R (15 points)}

A cohort study follows 500 smokers and 500 non-smokers for 10 years to assess heart disease incidence:

\begin{center}
\begin{tabular}{l|cc}
& \textbf{Heart Disease} & \textbf{No Heart Disease} \\
\hline
\textbf{Smoker} & 75 & 425 \\
\textbf{Non-smoker} & 35 & 465 \\
\end{tabular}
\end{center}

Use R to perform a chi-square test on this data (without continuity correction). What is the chi-square test statistic?

\begin{enumerate}[label=\Alph*.]
\item 8.234
\item 12.156
\item 16.343 \textbf{(CORRECT)}
\item 21.450
\end{enumerate}

\vspace{0.5cm}

\textbf{Expected R code solution:}
\begin{verbatim}
data <- matrix(c(75, 425, 35, 465), nrow=2, byrow=TRUE)
chisq.test(data, correct=FALSE)
\end{verbatim}

\textit{R output:}
\begin{verbatim}
	Pearson's Chi-squared test

data:  data
X-squared = 16.343, df = 1, p-value = 0.00005285
\end{verbatim}

\textit{Explanation:} The chi-square test statistic is $\chi^2 = 16.343$. With $df = 1$ and $p < 0.0001$, we reject $H_0$ and conclude there is very strong evidence of an association between smoking and heart disease.

\newpage

\section{Part C: Group Project Check-In (Completion Only, Ungraded)}

\subsection*{Progress Update Check-In}

\textbf{GRADING:} Completion only (any honest response receives credit)

Group leaders should submit a brief update (3-5 sentences) addressing:
\begin{enumerate}
\item Has your group met at least once since Week 7?
\item Have you loaded and explored your dataset in R?
\item Have you consulted at least one AI tool about your data?
\item Any blockers or questions?
\end{enumerate}

\vspace{12pt}

\section*{Points Summary}

\begin{itemize}
\item Part A (Multiple Choice): 70 points (7 × 10 pts) - AUTO-GRADED
\item Part B (R Computational): 30 points
  \begin{itemize}
  \item Question 8: Fisher's Exact Test (15 pts - AUTO-GRADED)
    \begin{itemize}
    \item Part A: R code submission (not graded - for reference only)
    \item Part B: p-value multiple choice (15 pts - AUTO-GRADED)
    \end{itemize}
  \item Question 9: Chi-Square Test (15 pts - AUTO-GRADED)
    \begin{itemize}
    \item Part A: R code submission (not graded - for reference only)
    \item Part B: test statistic multiple choice (15 pts - AUTO-GRADED)
    \end{itemize}
  \end{itemize}
\item Part C (Group Check-In): Completion only (MANUALLY GRADED)
\item \textbf{Total: 100 points (70 + 30 auto-graded, 0 manually graded except completion check)}
\end{itemize}

\end{document}
